\scanpage{92}
% archive sheet number: 91
% block text | lines 1-2
La suite exacte de coh. montre que $f$ est dans l'image de
% block display | line 3
\[
H^1(\Gamma,m^nA/m^{n+1}A)\longrightarrow H^1(\Gamma,A/m^{n+1}A^*)
\]
% block text | lines 4-5
i.e. est $\sim$ \inserted{dans $A/m^{n+1}A^*$} à un cocycle de la forme suivante \deleted{$\unsure{1+g\gamma|1+g}$}
% block display | line 6
\[
1+g(\unsure{\gamma|\sigma})
\]
% block text | line 7
où $g(\gamma)\in m^nA$ satisfaisant à
% block display | line 8
\[
g(\gamma\gamma')\equiv g(\gamma)+\gamma g(\gamma')\qquad(m^{n+1}A)
\]
% block text | lines 9-11
\deleted{Par hypothèse, on} On peut dans le \unsure{raisonnement} supposer $m^{n+1}=0$.

\deleted{Donc} L'hypothèse \deleted{\unsure{implique|indique}} implique qu'on a
% block display | line 12
\[
1+g(\gamma)=u^{-1}\gamma(u)\qquad\text{pour }\gamma\in\Lambda
\]
% block text | line 13
i.e.
% block display | line 13
\[
\gamma(u)-u=ug(\gamma)\deleted{\unsure{\in|\equiv|=}}\qquad\text{pour }\gamma\in\Lambda.\quad\deleted{\unsure{Comme}}
\]
% block insertion | above line 14
\inserted{\deleted{$mg(\gamma)=0$}}
% block text | lines 14-16, cancelled boxed block
\deleted{\unsure{Comme}}
% block display | line 14, cancelled boxed block
\[
\deleted{mg(\gamma)=0,\quad\text{i.e.}\quad ug(\gamma)=\unsure{\varepsilon|\epsilon}(u)g(\gamma)}
\]
% block text | line 15, cancelled boxed block
\deleted{i.e.}
% block display | line 15, cancelled boxed block
\[
\deleted{g(\gamma)=\gamma(\varepsilon(u)^{-1}u)-\varepsilon(u)^{-1}u}
\]
% block text | line 16, cancelled boxed block
\deleted{\unsure{on rappelle|en rappelant} que $\Gamma$ opère \uline{\unsure{trivialt|trivlt}} sur $O/m$ !)}

% block text | lines 17-23, cancelled block
\deleted{Cela signifie que $\gamma\mapsto ug(\gamma)$ est $\sim0$ dans $H^1(\Lambda,A)$ (additif !). \deleted{\unsure{Donc}} Or, $A$ est un $(\Gamma$-$O)$-module de longueur finie, on prouve par récurrence sur la longueur d'un tel module $M$ que l'application $H^1(\Gamma,M)\longrightarrow H^1(\Lambda,M)$ est \uline{injective}. \unsure{Or|Car}}

% block text | lines 24-25, doubly cancelled bottom lines
\deleted{\deleted{\unsure{Il suffit si} $M$ est \unsure{de longueur un|simple}, \unsure{ce qui est bien connu} : \illegible{1} par la \unsure{suite|raison} \illegible{1}}}

% block margin note | left margin, written bottom to top beside lines 1-24
% margin note written sideways
\margin{\deleted{Dans ce \inserted{dernier} cas, \deleted{\unsure{on a}} en effet $mM=0$ (sinon on aurait $mM=M$, d'où $M=0$ par Nakayama) donc $M$ est \deleted{\unsure{en}} en fait un vectoriel sur $k$ dans lequel $\Gamma$ opère par transformations \uline{$k$-linéaires}. On sait bien que $H^1(\Gamma,M)\longrightarrow H^1(\Lambda,M)$ est alors injectif.}}
